\setpart{框架能力 · 对照与缺口}

\begin{frame}[t]{能力对照（一）：输入与输出}
\scriptsize
\renewcommand{\arraystretch}{1.12}
\setlength{\tabcolsep}{3pt}
\begin{tabularx}{\textwidth}{>{\bfseries}p{2.9cm} Y Y Y}
\toprule
\rowcolor{Mist}
能力（Qwen3-Omni 语境） & \textcolor{Teal}{M*} & \textcolor{Orange}{vLLM-Omni} & \textcolor{Blue}{SGLang-Omni} \\
\midrule
输出流式（中间结果返回） & 有：NDJSON，每个 token 与音频块一条 & 有：文本与音频增量事件 & 有：SSE，音频 PCM\_F32 \\
运行中的外部输入（追加） & \textcolor{Red}{无}：Qwen3-Omni 未声明输入流，追加被拒 & 部分：帧与音频缓存，查询时使用 & 部分：会话输入，按轮次切分 \\
输入形式 & 文本；图像（PNG/JPEG/WebP/GIF）；音频与视频（随整段请求提交，采样率未核对） & 音频 PCM16 16 kHz；JPEG/PNG 帧；文本 & 音频 PCM16；摄像头帧；文本 \\
多模态时间语义 & 无 & 部分：帧带 frame\_id，返回 latest\_pts\_ms & 未核对 \\
阶段间流式（输出侧） & 有：\code{StreamingGraphEdge} & 有：async\_chunk & 有（阶段运行时，未逐项核对） \\
\bottomrule
\end{tabularx}
\par\takeaway{三个框架都有输出流式。输入流式只有 vLLM-Omni 与 SGLang-Omni 的会话式入口可用，M* 对 Qwen3-Omni 没有。}
\par\memo{依据见前面三页（M* 第 1--3 页，vLLM-Omni 第 1--4 页，SGLang-Omni 第 1 页）。未核对的格子不作结论。}
\end{frame}

\begin{frame}[t]{能力对照（二）：打断、取消与增量}
\scriptsize
\renewcommand{\arraystretch}{1.12}
\setlength{\tabcolsep}{3pt}
\begin{tabularx}{\textwidth}{>{\bfseries}p{2.9cm} Y Y Y}
\toprule
\rowcolor{Mist}
能力 & \textcolor{Teal}{M*} & \textcolor{Orange}{vLLM-Omni} & \textcolor{Blue}{SGLang-Omni} \\
\midrule
打断（barge-in） & \textcolor{Red}{无}：\code{response.cancel} 被拒绝 & 有：双工端点，Qwen3-Omni 插件 \code{supports\_barge\_in} & 有：会话 INTERRUPT，按轮次 \\
新查询 / 新输入对旧输出的处理 & 未见 & 抢占：取消旧查询，旧输出丢弃 & 截断或重问 \\
中止在途计算 & 部分：排队请求丢弃，已启动的一步跑完 & 有：取消 $\to$ 引擎 abort，步间生效（GPU 停止时机推断） & 部分：断开 $\to$ abort；是否立即停未核对 \\
增量 prefill（只算新增上下文） & \textcolor{Red}{无}（未找到） & 部分：需前缀缓存、只出文本，查询间隙预热 & 未核对（历史保留，\code{thinker\_session.py}） \\
推理中并入新输入（改变正在进行的生成） & \textcolor{Red}{无} & \textcolor{Red}{无} & \textcolor{Red}{无} \\
Qwen3-Omni 可用的入口 & \code{/generate}（窗口式，本仓库未测） & 聊天、视频会话、\code{/v1/realtime}、双工端点 & 聊天、\code{/v1/realtime}（需 \code{--enable-realtime}） \\
\bottomrule
\end{tabularx}
\par\takeaway{三个框架都没有"推理中并入新输入"。打断与取消是 vLLM-Omni 与 SGLang-Omni 已有的，M* 没有。}
\par\memo{依据见前面各页。"推理中并入新输入"这一行是三个框架共同的缺口：未找到任何使正在进行的生成吸收新输入的机制。}
\end{frame}

\begin{frame}[t]{要达到目标，各框架需要改哪一层}
\scriptsize
\renewcommand{\arraystretch}{1.12}
\setlength{\tabcolsep}{3pt}
\begin{tabularx}{\textwidth}{>{\bfseries}p{2.6cm} Y Y Y}
\toprule
\rowcolor{Mist}
目标能力 & M* & vLLM-Omni & SGLang-Omni \\
\midrule
窗口式提交（Qwen3-Omni） & 外部运行时即可（\code{/generate} 逐次提交） & 外部运行时即可（聊天接口） & 外部运行时即可（聊天接口） \\
会话式输入与查询 & 需要模型层改动：声明输入流，实现 \code{process\_stream\_chunk} & 现成（视频会话 API），需要的是音频输出时的增量 prefill & 现成（会话入口），需要验证其 GPU 层面的行为 \\
推理中并入新输入（核心缺口） & 调度与引擎层：把新帧、音频接入正在进行的序列，并处理位置与 KV & 同左，并加上增量编码（音频、视频帧） & 同左 \\
打断到 GPU 停止 & 取消路径要进入 worker 的步级检查 & 已有取消到引擎 abort，需验证打断延迟 & 需验证 abort 是否立即停止批次 \\
Pi0.5：同连接接收新观测（E1） & OpenPI 端点的接收与推理解耦（\code{openpi.py:285}；\code{openpi.rs:566}） & 服务层（\code{connection.\allowbreak{}py:228--278}） & 服务层（\code{ws\_utils.py}） \\
Pi0.5：推理中止（E2） & 去噪步之间加取消点 & 同左（框架内的去噪循环不经过取消路径） & 同左 \\
\bottomrule
\end{tabularx}
\par\takeaway{窗口式不需要改框架；会话式和打断在 vLLM-Omni 与 SGLang-Omni 已经存在；真正的缺口是推理中并入新输入，三个框架都要改引擎。}
\par\memo{改动层的划分：外部运行时（不改框架）→ 协议与服务层 → 调度层 → 模型层 → 引擎层。}
\end{frame}
